\subsection{as4の半単純化を用いる包絡候補の縮小}
\label{sec:as4-h3-reductions}
本節は\texttt{uedakazushi/as4}のコミット
\texttt{60eacafa28156a5fc62f74eae004d64a84675f4b}を参照する．
輸入するのは，\texttt{UNIPOTENT\_UNTWIST.tex}の一般帰着と，
\texttt{QUANTUM\_UNIPOTENT\_LIFTS.tex}の量子基底の全被覆，
および後述する二つのCY構造表示である\cite{AS4pinned}．
同repoの全CY統合定理・全AS4被覆を，独立監査済み入力として用いるものではない．
全て通常のgraded・1周期代数についての主張であり，
最小周期2・4のAS正則$\mathbb Z$代数の例をそのままH3へ移すことはしない．
H3の二次包絡は，指定構造対象と反Serre orbitから得たgraded環$R$の$Q(R)$である．
別のhelixの周期表示と比較する際にはこの対象の同定が必要である．

\paragraph{中山規約．}
as4は右巡回$\tau(abcd)=dabc$と
$(\nu\otimes1^{\otimes3})\tau(w)=-w$を用いる．
その量子基底では$\nu_Q(x_i)=(\prod_jq_{ji})x_i$であり，
本稿の$\mu_Q(x_i)=(\prod_jq_{ij})x_i$の逆である．
従って本節では$\nu=\mu^{-1}$へ変換する．
右twistの積は両方で$u*v=uH^{\deg u}(v)$である．
単冪$H$のhomological determinantは1なので，
$B=C^H$では
\[
 \mu_B=\mu_CH^{-4},\qquad \mu_C H=H\mu_C.
\]
この符号は独立にも確認できる．次数$n$で$H^{-n}$を作用させる写像は
右twist $C^H$を左twist $C^{\widetilde{H^{-1}}}$へ同型にする．
\cite[Theorem 0.3]{RRZ}の左twist公式へ$H^{-1}$を入れると上の式となる．
指数の符号を通常の行列式で置き換えない．

\begin{proposition}[H3包絡のスペクトル必要条件]
\label{prop:h3-spectral-filter}
$B$がH3の包絡で，$\mu_B|_{B_1}$の固有値を重複込みで
$\lambda_0,\ldots,\lambda_3$とする．ある$a\in\mathbb N^4$，$|a|=3$について
\[
 \lambda_0^{a_0}\lambda_1^{a_1}\lambda_2^{a_2}\lambda_3^{a_3}=1
\]
となる必要がある．さらに固有値が一種類$\lambda$だけなら$\lambda=1$である．
特に半単純部分が$-I$，$iI$，$-iI$である全Jordan層はH3包絡から除外される．
\end{proposition}
\begin{proof}
$\mu_B(G)=G$なので，$B_3$上に固有値1が必要である．
$B_3$は$B_1^{\otimes3}$の商であり，その固有値は三つの一次固有値の積に限られる．
これで最初の条件を得る．
固有値が$\lambda$一種類なら$\mu_B=\xi_\lambda U$と書け，$U$は単冪自己同型である．
三次作用の固有値は全て$\lambda^3$なので$\lambda^3=1$である．
またnoetherian Koszul AS正則環の中山はhomological determinant 1を持つ
\cite[Theorem 6.3]{RRZ}．単冪$U$の一次元Koszul topへの作用は1，
次数スカラー$\xi_\lambda$のhomological determinantは$\lambda^4$だから
$\lambda^4=1$となる．従って$\lambda=1$である．
一般の自己同型のhomological determinantを通常の$4\times4$行列式と同一視してはいない．
\end{proof}

\begin{lemma}\label{lem:h3-top-nonseparable}
H3包絡 $B$ の一次成分を非零部分空間の直和 $V=A\oplus C$ とする。
四次Koszul topを表す非零tensor $w\in V^{\otimes4}$ は
$A^{\otimes4}\oplus C^{\otimes4}$ に属さない。
\end{lemma}
\begin{proof}
Koszul双対のFrobenius積 $E_2\otimes E_2\to E_4$ は完全対合である。
従って $w$ の二次flatteningの像は六次元関係空間 $R$ 全体である。
もし $w=w_A+w_C$ なら、このflatteningはブロック対角であり、
$R=R_A\oplus R_C$、$R_A\subset A^{\otimes2}$、$R_C\subset C^{\otimes2}$ となる。
従って $B=T(A)/(R_A)*T(C)/(R_C)$ は二つの連結次数付き代数の自由積である。
$B$ は整域なので、各因子内の非零一次元 $a\in A,c\in C$ は
それぞれ多項式部分代数 $k[a],k[c]$ を生成する。
自由積の正規形により $k[a]*k[c]=k\langle a,c\rangle$ が $B$ に埋め込まれる。
その指数成長は $H_B(t)=(1-t)^{-4}$ に矛盾する。
\end{proof}

\begin{proposition}[二固有値の有限スペクトル制約]
\label{prop:h3-two-eigenvalue-filter}
H3包絡の $\mu_B|_{B_1}$ が相異なる二固有値 $\alpha,\beta$ を持つなら、
その半単純部分の位数は $2,3,5,6,9$ のいずれかである。
Jordanの単冪部分の消滅は、この結論には含めない。
一つの原始90乗根 $\zeta$ を固定したとき、可能な非順序固有値集合は
次の16組に限られる：
\[
\begin{array}{c|l|c}
\text{位数}&\{\alpha,\beta\}&\text{組数}\\\hline
2&\{1,-1\}&1\\
3&\mu_3\text{ の相異なる二元}&3\\
5&\{\xi^k,\xi^{2k}\},\ \xi\text{ 原始5乗根},\ k\in(\mathbb Z/5)^*&4\\
6&\{\xi^k,\xi^{2k}\},\ \xi\text{ 原始6乗根},\ k=1,5&2\\
9&\{\xi^k,\xi^{-3k}\},\ \xi\text{ 原始9乗根},\ k\in(\mathbb Z/9)^*&6
\end{array}
\]
各組のAS4存在・H3許容性を主張するものではない。
\end{proposition}
\begin{proof}
$A,C$ を $\alpha,\beta$ に対する一般化固有空間とする。
$\mu_B(G)=G$ と $B_3$ が $V^{\otimes3}$ の商であることから、
ある $0\le a\le3$ について $\alpha^a\beta^{3-a}=1$ である。
一方、\cite[Theorem 6.3]{RRZ}より最高Koszul tensorは $\mu_B^{\otimes4}$ に固定される。
補題\ref{lem:h3-top-nonseparable}から、$w$ には両一般化固有空間を含む成分があり、
その $A$ 因子の個数を $1\le c\le3$ とすると
$\alpha^c\beta^{4-c}=1$ となる。
一般化固有空間でも、このtensor成分の唯一の固有値はこの積であるから、
非零固定成分の存在から同じ等式が従う。

二つの等式を消去すると
\[
 \alpha^{4a-3c}=\beta^{4a-3c}=1.
\]
$4a-3c$ は零にならず、その絶対値の可能性は $1,2,3,5,6,9$ である。
絶対値1では両固有値が1となり、相異なるという仮定に反する。
従って二固有値の位数の最小公倍数は $2,3,5,6,9$ のいずれかとなる。
$\mu_{90}$ 内で二つの重み等式を解くと表の16組を得る。
この最後の有限整数演算は
\texttt{verify\_\allowbreak{}cubic\_\allowbreak{}two\_\allowbreak{}eigenvalue\_\allowbreak{}filter.py}
で再現し、tensorと自由積の議論は上の証明に依存する。
\end{proof}

\begin{lemma}\label{lem:h3-bigraded-top-index}
$B$ が $V=A\oplus C$ の個数による二重次数を持ち、
四次Koszul topの二重次数が $(c,4-c)$ であるとする。
$p=\dim A$, $q=\dim C$ とすると $c=p$, $4-c=q$ であり、
\[
 H_B(s,t)=\frac1{(1-s)^p(1-t)^q}.
\]
\end{lemma}
\begin{proof}
二重次数付きKoszul解消から $H_B(s,t)=1/P(s,t)$ を得る。
$P$ は全次数4以下で、最高次数部分はただ $s^ct^{4-c}$、
一次部分は $-ps-qt$ である。
任意の正整数 $m,n$ について、正の再次数付けは元の次数と上下に比較できるので、
$H_B(u^m,u^n)$ は $u\to1^-$ で正確に4位の極を持つ。
従って $P$ の $(1,1)$ におけるTaylor展開の4次未満の部分は全て零である。
実際、最小非零次数が4未満なら、その斉次部分を $(m,n)$ で評価した値が
全ての正整数対で零となり、その斉次多項式自体が零となって矛盾する。
最高次数項は一つなので
$P(s,t)=(s-1)^c(t-1)^{4-c}$ となる。
一次係数の比較で $p=c,q=4-c$ を得る。
\end{proof}

\begin{lemma}\label{lem:h3-binary-cubic-factors}
二生成元・一二次関係の次数付き整域では、全ての非零斉次三次式は一次式の積に分解する。
\end{lemma}
\begin{proof}
関係の $2\times2$ 係数行列は可逆である。
階数一なら関係自体が二一次式の積であり、整域性に反するためである。
代数閉体上で関係の可換対称部分の等方方向を選べば、基底を変えて
\[
 yx=qxy+\kappa y^2,\qquad q\ne0
\]
と書ける。これはPBW表示であり、$y$ は正規正則である。
$y$ を逆にして $z=xy^{-1}$ と置くと $yz=(qz+\kappa)y$ となる。
次数 $d$ の式は $f=p(z)y^d$, $\deg p\le d$ と一意に書ける。
$p$ を一次式へ分解し、その根を $a_1,\ldots,a_s$ とすると、
\[
 b_j=q^{j-1}a_j+\kappa\sum_{i=0}^{j-2}q^i
\]
によって、非零定数倍を除き
$f=(x-b_1y)\cdots(x-b_sy)y^{d-s}$ となる。
これは局所化での等式だが、両辺は元の整域内にあるので元の環でも成立する。
$d=3$ では必ず低次数の非零因子に分解する。
\end{proof}

\begin{theorem}[二固有値H3包絡の中山半単純部分]
\label{thm:h3-two-eigenvalue-refined}
H3包絡の中山自己同型が相異なる二固有値を持つなら、その半単純部分の位数は
2または3である。従って固有値集合は
\[
 \{1,-1\},\quad \{1,\omega\},\quad\{1,\omega^2\},\quad\{\omega,\omega^2\},
 \qquad \omega^2+\omega+1=0
\]
に限られる。重複度、Jordan型、存在およびH3許容性は別の判定である。
\end{theorem}
\begin{proof}
命題\ref{prop:h3-two-eigenvalue-filter}の残る位数 $5,6,9$ を排除する。
その12個の固有値集合では、比 $\alpha/\beta$ の位数もそれぞれ $5,6,9$ である。
従って4次topの重み等式を満たす $\alpha$ の個数 $c$ は一意であり、
$w$ は二重斉次となる。その二次flatteningの像 $R$ も二重斉次なので、
$B$ は二重次数を持つ。補題\ref{lem:h3-bigraded-top-index}から
一般化固有空間の次元は $(c,4-c)$ である。
同様に三次関係 $G$ の $\alpha$ 個数 $a$ も一意となる。

位数5では $(a,c)=(1,3)$ または $(2,1)$ だけである。
次元分解が $3+1$ のとき、二重Hilbert式から混合二次関係は三本である。
一次元側の生成元を $z$ とすると、混合関係の
$z\otimes A$ と $A\otimes z$ への両射影は整域性により単射である。
従って両方とも同型となり、$za=\sigma(a)z$、$\sigma$ は可逆な一次変換となる。
三次 $G$ の二重次数は、一次元側を二個含む型なので、
正規順序化により $G=\ell z^2$ または $G=z^2\ell$ と分解する。
これは $B/(G)$ の整域性に反する。

位数9では $(a,c)=(0,3)$ または $(3,1)$ だけであり、
$G$ は一次元側の純三乗となるので同様に不可能である。

位数6では $c=2$、$a=0$ または3となる。
純粋な各一般化固有空間の二次関係は一本である。
非負の二重次数により、その二変数部分代数にはこの一本以外の関係は生じず、整域である。
$G$ は片方の部分代数の三次式だけからなるので、
補題\ref{lem:h3-binary-cubic-factors}により分解し、再び整域性に反する。
以上で12組を全て排除する。
\end{proof}
この定理は全H3に対する必要条件であり、as4の未監査な全スペクトルatlasを輸入しない。
三種類以上の固有値を持つ層をこの定理で排除することも、
残る単冪部分を自動的に零とすることもない。

\begin{lemma}\label{lem:h3-binary-normalizer-mu3}
二生成元・一二次関係の次数付き整域 $T$ に非零正規三次 $F$ があり、
その正規化自己同型 $\theta|_{T_1}$ の固有値が原始三乗根 $\rho$ 一種類だけであるとする。
座標と非零係数を取れば
\[
 T=k\langle x,y\rangle/(xy-\rho^{-1}yx),\qquad
 F=x^2y,\qquad \theta=\rho I
\]
となる。
\end{lemma}
\begin{proof}
補題\ref{lem:h3-binary-cubic-factors}の二変数表示を量子平面またはJordan平面に直す。
可換平面では $F$ の正規化は恒等なので不可能である。
Jordan平面 $yx=xy+y^2$ の次数付き自己同型は
$x\mapsto ax+by$, $y\mapsto ay$ である。
これは一般の二次元行列を関係へ代入して係数を比較すれば従う。
$F=p(z)y^3$, $z=xy^{-1}$ と書くと
$Fy=\rho yF$ は $p(z)=\rho p(z+1)$ となる。
最高係数から $\rho=1$ となって矛盾する。

量子平面 $xy=qyx$, $q\ne1$ の次数付き自己同型は対角または反対角である。
実際、一般変換の関係の $x^2,y^2$ 係数は
$(1-q)ac,(1-q)bd$ なので、可逆性と合わせればよい。
反対角可逆行列は相異なる二固有値を持つため、$\theta=\rho I$ となる。
PBW単項式 $x^uy^{3-u}$ の正規化指標は $x,y$ 上で
$q^{u-3},q^u$ である。
両方が $\rho$ となるには $q^3=1$ であり、$u$ は1または2の一つに限られる。
必要なら $x,y$ を交換して $F=x^2y$, $q=\rho^{-1}$ とする。
\end{proof}

\begin{theorem}[二固有値・位数3の二構造]
\label{thm:h3-order3-two-blocks}
H3包絡の中山半単純部分が位数3で、固有値が二種類だけであるとする。
次のいずれかに限られる。
\begin{enumerate}
\item 固有値集合が $\{1,\rho\}$、$\rho$ は原始三乗根なら、
固有値1の一般化固有空間は直線 $kz$、$\rho$ の空間は三次元である。
$z$ は正規正則で $\mu_B(z)=z$、
\[
 B=S[z;\sigma],\qquad G=F+c z^3,
 \quad 0\ne F\in S_3,\quad \sigma(F)=F
\]
と書ける。$S$ は三生成元noetherian quadratic Koszul AS3であり、
$F$ は $S$ 内でも正規正則、その正規化は $\mu_B^{-1}|_S$ である。
実際には $c\ne0$ が必要であり、$z$ の定数倍で $G=F+z^3$ と規格化できる。
さらに
\[
 \sigma^3=\mu_B^{-1}|_S,\qquad \mu_S=\sigma^{-2}
\]
が必要である。
\item 固有値集合が $\{\omega,\omega^2\}$ なら、両重複度は二であり、
$\mu_B$ 自体が半単純である。座標を取って
\[
 \begin{gathered}
 \mu_B=\operatorname{diag}(\omega,\omega,\omega^2,\omega^2),\qquad
 xy=\omega yx,\quad zt=\omega^2tz,\\
 G=x^2y+z^2t
 \end{gathered}
\]
となる。残る四関係は $\langle x,y\rangle$ と $\langle z,t\rangle$ の混合関係で、
両順序の混合積の空間への射影は同型である。
従って $x^iy^jz^kt^l$ は $B$ のPBW基底となる。
\end{enumerate}
これは残る関係係数のAS4存在・H3許容性の十分条件ではない。
\end{theorem}
\begin{proof}
二固有値の比の位数は3なので、$V^{\otimes2}$ の三つの一般化固有値
$\alpha^2,\alpha\beta,\beta^2$ は相異なる。
$R$ はそれらの一般化固有空間に分解し、$B$ は二重次数を持つ。
一次元のKoszul topも二重斉次となる。
混合topの重みと補題\ref{lem:h3-bigraded-top-index}により、
重複度は第一の場合 $1+3$、第二の場合 $2+2$ に限られる。
三次の固定重みは両方とも純粋な三次二成分だけなので、$G=F_A+F_C$ と書ける。

第一の場合、二重Hilbert式から混合関係は三本である。
一次元側を $kz$ とすると、両射影の核に非零元があれば
$a z=0$ または $z a=0$ となり、整域性に反する。
従って混合関係は可逆な写像 $\sigma$ のgraphであり、$za=\sigma(a)z$ となる。
純三変数部分代数を $S$ とすると、この書換えとHilbert式から実際に $B=S[z;\sigma]$ である。
$S\simeq B/(z)$ はnoetherianである。
二重次数付きKoszul解消の $z$ 次数零部分は $S$ 上の有限線型解消を与え、
正規一次元に対するRees lemmaと合わせて $S$ はAS3となる。
$\mu_B(z)=z$ なので $Gz=zG$、従って $\sigma(F)=F$ である。
$F=0$ は純三乗切断となるため排除される。
$Gv=\mu_B^{-1}(v)G$ の二重次数ごとの比較から $F$ の正規式を得る。
$c=0$ なら $G=F\in S_3$ であり、その正規化
$\theta=\mu_B^{-1}$ は $S$ を保ち $z$ を固定する。
補題\ref{lem:cubic-skew-cone-vertex}の正規版により、
$\pi k[z]$ の自己Extが全次数で非零となってH3に反する。
従って $c\ne0$ であり、$z^3$ の係数比較が
$\sigma^3=\mu_B^{-1}|_S$ を与える。
代数閉性により $z$ を定数倍すれば $c=1$ となる。
正規一次元 $z$ の正規化は $S$ 上で $\sigma$ なので、
正規商の中山公式\cite[Lemma 1.5]{RRZ}は
$\mu_S=\mu_B|_S\circ\sigma=\sigma^{-2}$ を与える。

第二の場合、純粋な各二変数部分代数は一本の二次関係で定義される。
二重次数が非負なので、混合関係から新しい純粋関係は生じない。
$F_A,F_C$ の一方が零なら、他方は補題\ref{lem:h3-binary-cubic-factors}により分解するため、
両方とも非零である。
正規式を二重次数で分けると、それぞれの正規化は
$\mu_B^{-1}$ の制限であり、その唯一の固有値は $\omega^2,\omega$ である。
補題\ref{lem:h3-binary-normalizer-mu3}により二つの量子平面と表示した三次式を得る。
特に各二変数空間上の中山作用はスカラーとなり、単冪部分は消える。

最後に混合関係の射影を確認する。
$w$ の純二次二ブロックは、Frobenius対合から非零定数倍の
$(xy-\omega yx)\otimes(zt-\omega^2tz)$ である。
巡回式により、混合flatteningの $AC$ から $CA$ へのブロックは、
これら二つの可逆な $2\times2$ 係数行列のKronecker積に、
非零スカラーと置換を施したものである。従って可逆である。
逆方向も同様であり、混合flatteningの全階数は4なので、
その像 $R_{1,1}$ の両射影は同型となる。
これで混合積を常に $AC$ 順へ移せる。
得られた標準単項式は全語を張り、その個数は $H_B=(1-u)^{-4}$ に一致するため基底となる。
\end{proof}
重み、Jordan二変数の最高係数比較、量子二変数の正常化指標、
混合flatteningの可逆ブロックは
\texttt{verify\_\allowbreak{}cubic\_\allowbreak{}order3\_\allowbreak{}two\_\allowbreak{}eigenvalue\_\allowbreak{}reduction.py}
で有限恒等式を独立に確認する。
混合係数の必要系は次の定理で全て判定する。第一の場合の $c=0$ の一般正規cone障害は別に扱う。

\begin{theorem}[二つの非自明三乗根固有値の全排除]
\label{thm:h3-order3-two-nontrivial-excluded}
H3包絡の中山固有値集合は $\{\omega,\omega^2\}$ にならない。
従って二固有値の場合に残る集合は
$\{1,-1\},\{1,\omega\},\{1,\omega^2\}$ の三つだけである。
\end{theorem}
\begin{proof}
定理\ref{thm:h3-order3-two-blocks}の第二の場合を仮定する。
$A=\langle x,y\rangle$, $C=\langle z,t\rangle$ とし、
$w$ の $AA\,CC$ 成分を $(xy-\omega yx)\otimes(zt-\omega^2tz)$ に規格化する。
その係数はFrobenius対合により非零なので、この規格化は可能である。
巡回式で非交互の語の係数が全て決まり、交互部分は
$AC=(xz,xt,yz,yt)$ 順の対称行列
\[
 H=\begin{pmatrix}
 h_0&h_1&h_2&h_3\\h_1&h_4&h_5&h_6\\
 h_2&h_5&h_7&h_8\\h_3&h_6&h_8&h_9
 \end{pmatrix}
\]
だけで記述される。対称性は二回の巡回による指標が $\omega\omega^2=1$ であるためである。
混合flatteningのオフ対角は可逆なので、
階数4条件は16個のSchur成分の消滅である。
さらに四つの重なり $zyx,tyx,tzx,tzy$ のPBW条件と、
\[
 (z^2t)a=\omega^2a(z^2t)\ (a=x,y),\qquad
 (x^2y)c=\omega c(x^2y)\ (c=z,t)
\]
を正規順序化する。零と重複を除くと68個の必要多項式 $I_i$ を得る。

保存した多項式証明書は $I X=\Gamma$ という30個の恒等式であり、
元の68式を再生成して全て独立に検証した。
従ってGr\"obner基底の正しさを仮定せず、$I=0$ から $\Gamma=0$ が従う。
これらの式の一つは
\[
 (h_5-\omega^2)(3h_5-\omega-2)=0
\]
である。
第一の根では $h_3=-\omega^2$, $h_5=\omega^2$、他の係数は全て零となる。
これはskew係数
\[
 (q_{xy},q_{xz},q_{xt},q_{yz},q_{yt},q_{zt})
 =(\omega,\omega,\omega^2,\omega^2,1,\omega^2)
\]
の環である。定理\ref{thm:skew-finite-37}の必要リストでは、
三乗根のskew H3包絡は全てCYであるから、この非CY包絡は排除される。

第二の根では、$p,q\ne0$ を用いて全解が
\[
\begin{gathered}
 h_0=2(\omega+2)pq,\quad h_1=p,\quad h_2=q,\quad
 h_3=-(2\omega+1)/3,\quad h_5=(\omega+2)/3,\\
 h_4=(\omega-1)p/(3q),\quad h_7=(\omega-1)q/(3p),\\
 h_6=-\omega/(3q),\quad h_8=-\omega/(3p),\quad
 h_9=(2\omega+1)/(9pq)
\end{gathered}
\]
となる。必要性には、$h_0h_9=2\omega/3$、
$h_1h_8=h_2h_6=-\omega/3$、
$6h_1h_2+(\omega-1)h_0=0$ と、
$2h_1^2+h_0h_4=2h_2^2+h_0h_7=0$ を使う。
逆代入も68式を全て満たす。
$u^3=q/p$ を選び、$\operatorname{diag}(u,up,1,q)$ で生成元を変えると、
$G$ の直線を保ったまま $p=q=1$ となる。
従ってこの枝では一つの候補だけを調べればよい。

この候補で $\tau^3=\omega$ を選ぶと、次の文字どおりの積恒等式が成立する：
\[
\begin{aligned}
 \tau(x^2y+z^2t)=(x+\tau z)\bigl[&\tau xy+zt\\
 &+\tau^2\{\omega^2xz-\tfrac{2\omega+1}{3}(xt+yz)-\tfrac13yt\}\bigr].
\end{aligned}
\]
三次の全PBW係数を $\omega^2+\omega+1=0$, $\tau^3=\omega$ のもとで独立に確認した。
両因子は次数1と2の非零元であり、$B/(G)$ は整域でない。
これで第二の枝も全て排除される。
この排除には候補のAS4存在十分性を示す必要はない。
\end{proof}
全68式、30個の多項式所属恒等式、二枝のパラメータ化、尺度変更、因子式は
\texttt{verify\_\allowbreak{}cubic\_\allowbreak{}order3\_\allowbreak{}two\_\allowbreak{}block\_\allowbreak{}incidence.py}
と小さいJSON証明書で再現する。
生成にはSingular 4.4.1を用いたが、保存した証明書の再生はSymPyだけで行える。

\begin{lemma}[単冪twist下の反標準三次関係の輸送]
\label{lem:h3-unipotent-transport}
$C$は左右noetherian quadratic Koszul AS4型14641で，$\mu_C$は半単純とする．
$H$は$\mu_C$と可換な単冪graded自己同型，$B=C^H$とする．
両者の次数付きベクトル空間を同定すれば
\[
 W_1(B)=\{F\in C_3:H(F)=F,\quad
                 Fv=H\mu_C^{-1}(v)F\ \text{for all }v\in C_1\}.
\]
従ってH3包絡の三次関係は，半単純化後の$C$でも正規で，
その正規化自己同型は$H\mu_C^{-1}$となる．
\end{lemma}
\begin{proof}
$0\ne F\in W_1(B)$なら定理\ref{thm:nakayama-cubic-linearization}の逆向きにより
$F$は正規正則で$\mu_B(F)=F$である．
$C_3$上の$\mu_C$は半単純，$H^{-4}$は単冪で，両者は可換である．
従って$\mu_CH^{-4}(F)=F$からそれぞれのJordan部分が$F$を固定し，
特に$H(F)=F$となる．
$B$での正規式を$C$の積に戻すと
\[
 F H^3(v)=\mu_B^{-1}(v)H(F)=H^4\mu_C^{-1}(v)F.
\]
$v$を$H^{-3}(v)$に置き換えれば表示の式を得る．
逆に二条件を仮定して同じ計算を逆向きに行うと
$F*v=\mu_B^{-1}(v)*F$となるので$F\in W_1(B)$である．
\end{proof}

\begin{theorem}[量子半単純化を持つH3包絡には非自明な単冪がない]
\label{thm:h3-quantum-unipotent-exclusion}
as4の標準半単純化が量子多項式環$Q(q)$になるH3包絡$B$は，
$Q(q)$自身とgraded同型である．非自明な単冪持上げには$W_1(B)=0$である．
従ってこの\emph{非skew持上げを含む部分全体}のH3幾何学的分類は
定理\ref{thm:all-skew-H3-geometric}の31パラメータ・五型と一致する．
\end{theorem}
\begin{proof}
as4の量子基底全被覆定理により$B=Q(q)^H$と書ける．
前補題で得た$Q(q)$内の非零正規三次元$F$の正規化自己同型は
$H\mu_Q^{-1}$である．
定理\ref{thm:skew-normal-cubics}のNewton多面体の証明は，
量子多項式環の正規三次元の正規化自己同型が座標ごとに対角となることを示した．
$\mu_Q$も対角なので$H$は対角となる．単冪かつ対角な$H$は恒等である．
逆に恒等持上げでは既存skew分類を適用できる．
この証明はGL4内のJordan型だけの分類や，Zhang同値とgraded同型の同一視を用いない．
\end{proof}

\begin{lemma}[H3の反標準環は整域]
\label{lem:h3-orbit-domain}
H3対象の$R$は整域である．
\end{lemma}
\begin{proof}
構造対象とその反Serre shiftはrank-one torsion-freeである．
非零射$\OO\to\tau^n\OO$はgeneric商で非零な単純対象間の射なので同型となり，
核はtorsionであるため零である．generic像が零なら像がtorsionとなり，
標的のtorsion-free性から元の射も零となるので，この可能性はない．
従って非零斉次切断の積は非零な単射の合成となる．
次数付き環の最高非零次数を取れば非斉次元についても零因子がない．
\end{proof}

\begin{theorem}[中心二平面Heisenberg基底の全単冪持上げの除外]
\label{thm:h3-central-quadratic-unipotent-exclusion}
$0\ne q\in k[b,c]_2$とし
\[
 C_q=k\langle b,c,u,v\rangle/(b,c\text{ central},\ [u,v]=q(b,c))
\]
とする．標準半単純化が$C_q$であるAS4代数はH3包絡にならない．
従ってas4の$bc$型・$b^2$型の全単冪持上げを，一般値だけでなく全特殊値で除外できる．
\end{theorem}
\begin{proof}
$C_q$はCYで，$k[b,c]$上PBW自由である．
その中心体$K=k(b,c)$へ局所化すると，$[u,v]=q\in K^*$なので
一次Weyl代数$A_1(K)$を得る．これは単純環であり，単元は$K^*$に限る．
非零正規元$F\in C_q$の局所化は非零正規元なので単元となり，$F\in K$である．
PBW基底から$C_q\cap K=k[b,c]$である．従って全ての正規元は中心的である．

$B=C_q^H$がH3包絡なら，輸送補題によって$C_q$上の正規化は$H$となる．
しかし正規元が中心的なので$H=1$である．
残る$G\in k[b,c]_3$は代数閉体上で一次式と二次式の積に分かれる．
それぞれの因子は三次商で非零なので$C_q/(G)$には零因子があり，
補題\ref{lem:h3-orbit-domain}に反する．
$q=0$はWeyl局所化の仮定外であり，可換包絡として保持する．
\end{proof}

\subsection{可換三変数の非零導分Ore包絡に対する全パラメータ障害}
\label{sec:ore-point-obstruction}
以下は一般点での中心三次消滅より強い障害であり，中心三次が存在する特殊値も扱う．

\begin{theorem}[Ore軸のpoint対象の無限自己Ext]
\label{thm:h3-polynomial-ore-exclusion}
$D=k[x_1,x_2,x_3]$，$0\ne\delta:D\to D$は次数1の導分，
$B=D[T;\delta]$，$\deg T=1$とする．
$B$がH3設定のnoetherian quadratic Koszul AS4条件を満たすとする．
すると$B$はH3包絡にならない．
より正確に，任意の$0\ne G\in W_1(B)$に対し$R=B/(G)$と置くと，
$M=R/(x_1,x_2,x_3)R\simeq k[T]$の像$P=\pi M$は非零で
\[
 \dim\Ext^n_{\qgr R}(P,P)=
 \begin{cases}1&n=0,\\3&n=1,\\4&n\ge2\end{cases}
\]
となる．特に$\gldim\qgr R=\infty$である．
\end{theorem}
\begin{proof}
Liu--Wang--Wu\cite[Theorem 3.3]{LWWOre}を$\sigma=1$，$D$の中山恒等へ適用すると
$\mu_B|_D=1$である．従って$G\in W_1(B)$は全ての$x_i$と可換である．
$G=\sum_{j=0}^m f_jT^j$，$f_m\ne0$と書くと，$m>0$のとき
ある$\delta(x_i)\ne0$に対する交換子の最高$T$次数係数は
$m f_m\delta(x_i)$となり，零ではない．よって$m=0$，$G\in D_3$である．
$G$の正規性から$\delta(G)\in(G)D$なので，
\[
 C=D/(G),\qquad R=C[T;\bar\delta].
\]
$J=(x_1,x_2,x_3)$は二側イデアルで，$M=R/J\simeq k[T]$となる．

まず$P$に対するline objectの高次Ext消滅を確かめる．
$B$は$D$上左右自由であり，$M$の$B$上の解消は
$D$の三変数Koszul解消を$B$へ誘導したものである．
従って$\underline\Ext_B^j(M,B)=0$ ($j\ne3$)である．
正規正則$G$の二項解消とderived Homのchange of ringsから，
$\underline\Ext_R^j(M,R)=0$ ($j\ne2$)となる．
正規化自己同型によるtwistはこの消滅次数を変えず，$J$も保つ．
$R$はAS Gorenstein次元3・parameter 1であり，局所双対から
$H^i_{R_+}(M)=0$ ($i\ne1$)を得る．従って商圏コホモロジーとの比較により
\[
 \Ext^j_{\qgr R}(\pi R(a),P)=0\quad(j>0,\ a\in\mathbb Z).
\]
$T$倍写像$M\to M(1)$は$R$線型で有限次元cokernelを持つので$P\simeq P(1)$である．
さらにsection比較は十分大きい次数で$\dim M_n=1$と一致するため
$\Hom(\pi R(a),P)=k$が全$a$で成立する．

次に$C$上の$k=C/(x_1,x_2,x_3)$の解消を取る．
$G=\sum_i x_iG_i$，$G_i\in D_2$と書き，Koszul複体に
homological次数2・内部次数3の変数$z$を加えて
\[
 d(e_i)=x_i,\qquad d(z)=\sum_iG_ie_i
\]
とするTate複体を用いる．$d^2(z)=G=0$である．
元のKoszul複体のhomologyは$k$と，この一次cycleが張る$k(-3)$だけである．
$z$次数によるspectral sequenceで$d(z^m)=m z^{m-1}d(z)$となり，
標数0により後者のhomologyを全て殺すので，これは$k$の自由解消となる．
行列検算では同値なdivided-power基底$z^{[m]}=z^m/m!$を用いる．
全微分成分が$(x_1,x_2,x_3)$に入るため最小であり，Betti数は
\[
 \sum_{n\ge0}\beta_nt^n=\frac{(1+t)^3}{1-t^2}
 =1+3t+4t^2+4t^3+\cdots.
\]

$R$は左$C$加群として自由なので，この解消を$-\otimes_CR$で誘導すると
$M$の有限階数graded自由$R$解消が得られる．
$\qgr$へ移した解消の各項は$P$に対して上で示したHom-acyclic対象である．
従ってfirst-quadrantのExt spectral sequenceで$\Ext^n(P,P)$をこのHom複体から計算できる．
全微分成分は$J$に入り$P$に零作用するので，Hom複体の微分は零である．
各自由summandから$k$が一つ得られ，表示の次元が従う．
これはchartへの未証明のOre局所化や，line objectの射影性を仮定した議論ではない．
\end{proof}

\begin{lemma}[可換三平面を保つCY Ore基底の全単冪持上げ]
\label{lem:h3-ore-unipotent-exclusion}
$C=D[T;\delta]$，$D=k[x_1,x_2,x_3]$，$\delta\ne0$，
$\operatorname{div}\delta=0$とし，$C$はCY AS4とする．
$D$を保つ任意の単冪graded自己同型$H$について，$C^H$はH3包絡にならない．
\end{lemma}
\begin{proof}
H3関係が存在すれば，輸送補題から$C$上の正規三次元$F$の正規化は$H$である．
$H(D)=D$なので$D\setminus0$のOre局所化へこの正規式を延ばせる．
局所化は$K[T;\delta]$，$K=k(x_1,x_2,x_3)$である．
$\delta\ne0$，標数0のときこの微分作用素環は単純である：
非零二側イデアルに最小$T$次数の元を取り，正次数なら
$\delta(a)\ne0$となる$a\in K$との交換子で次数を下げられるので矛盾する．
次数0の非零元は$K$で可逆である．
従って$F$は局所化で単元となる．単元の$T$次数は0なので$F\in K\cap C=D$である．
$D$は可換で$F\ne0$だから，$Fa=H(a)F$から$H|_D=1$を得る．
単冪性により$H(T)=T+h$，$h\in D_1$となる．

右twistの実際の交換子を計算すると，$C^H$も可換三変数上のOre拡大であり，
その導分は$\delta'=\delta-hE$，$E=\sum_i x_i\partial_i$となる．
生成元上の全六関係が一致し，両環のPBW Hilbert級数も一致するので，これは環の同型である．
もし$\delta'=0$なら$\delta=hE$であり，
$\operatorname{div}(hE)=4h$と$\operatorname{div}\delta=0$から$h=0$，$\delta=0$となって矛盾する．
従って$\delta'\ne0$であり，定理\ref{thm:h3-polynomial-ore-exclusion}を適用できる．
\end{proof}

\begin{corollary}[CY基底の四・五交換子層と三交換子beta層の全持上げの除外]
\label{cor:h3-cy-commutator-reduction}
H3包絡$B$の標準半単純化$C$がCYで
$c(C)=\dim(R_2(C)\cap\bigwedge^2C_1)\ge4$なら，$B$は可換多項式環である．
また標準半単純化がas4の三交換子$\beta$型となるH3包絡も存在しない．
元の$B$ではなく，半単純化$C$の交換子次元を用いていることに注意する．
特にPymの$S(2,3)$標準Ore族は，非零導分の\emph{全パラメータ}で除外される．
\end{corollary}
\begin{proof}
as4の\texttt{SCALAR\_LARGE\_COMMUTATORS.tex}の
\texttt{thm:scalar-large-commutators}は$c(C)\ge5$を可換型と
$[u,v]=b^2,bc$の二型へ帰着する．可換基底の非自明持上げは
定理\ref{thm:h3-quantum-unipotent-exclusion}で除外され，
後二者の全持上げは定理\ref{thm:h3-central-quadratic-unipotent-exclusion}で除外される．
$c(C)=4$は\texttt{CY\_FOUR\_COMMUTATORS.tex}の
\texttt{thm:cy-four-commutators-structure}により
$k[a,b,c][u;\delta_H]$，$\delta_Ha=0,\delta_Hb=H_c,\delta_Hc=-H_b$で，
$H_b,H_c$は独立だから$\delta_H\ne0$である．
同定される全自己同型は可換三平面$\langle a,b,c\rangle$を保つため，
前補題で全単冪持上げを除外できる．
同じく\texttt{CY\_THREE\_COMMUTATORS\_REDUCTION.tex}の
\texttt{thm:cy-three-beta}は$\beta$型を非零divergence-free導分のOre拡大へ帰着し，
全自己同型が同じ可換三平面を保つことを与える．前補題を適用する．
これらの輸入はgraded AS4の指定構造表示だけであり，
全CY atlasの他の層や全Z-algebraの被覆を独立検証したとするものではない．
\end{proof}
